\clearpage\begingroup\SourceLeaf{97b}\HistoricalHead{92}{L}{PHYSIK. BRIEFWECHSEL.}
\begin{center}[12.]\end{center}
\begin{center}\textsc{Gauss} an \textsc{Encke}.\quad Göttingen, 8. August 1834.\end{center}
\noindent\hspace*{\parindent}Beigehend übersende ich Ihnen, mein theuerster Freund, die am 6. und
7. d. hier gemachten Magnetischen Beobachtungen, die für sich selbst sprechen
werden, und von denen Sie beliebigen Gebrauch machen können. Mir selbst
hat mein leider noch fortdauerndes Uebelbefinden eine unmittelbare Theilnahme
diesmahl nicht erlaubt, ich habe mich darauf beschränkt, zuweilen einen kleinen
Satz genau gleichzeitig correspondirender Beobachtungen in der Sternwarte
mit der 25 pfündigen Nadel zu machen, welche, so wie ähnliche an früheren
Tagen, eine höchst merkwürdige Harmonie zeigen und beweisen, dass auch
die kleinen Sprünge innerhalb weniger Minuten, die uns früher öfters beunruhigt
haben, durchaus reell sind, obwohl erst künftig Vergleichung mit
scharfen entferntern Beobachtungen ausweisen wird, ob auch diese \textls[120]{sehr weit}
entlegene Ursachen haben. Jedenfalls sind die Intervalle von 5 zu 5 Minuten
keinesweges zu enge, sondern, zu Zeiten starker oder sprudelnder Anomalien
eher noch zu gross. Allein ohne ein sehr zahlreiches Personal würde es nicht
wohl thunlich sein, 44 Stunden hindurch jede Minute aufzuzeichnen, allein
wenn die scharfen Apparate erst noch mehr verbreitet sein werden, möchte
es nützlich sein, zuweilen solche correspondirende Beobachtungen mit sehr
engen Intervallen für eine kleine Anzahl von Stunden zu verabreden. Wir
haben hier schon zuweilen solche Aufzeichnungen gleichzeitig an beiden Apparaten,
5 bis 6 Resultate in der Minute, gemacht; versuchsweise sogar einmal
\textls[120]{dadurch} die Uhren verglichen, was sich hinterdrein auf 12 Sekunden richtig
fand. Ich bin geneigt zu glauben, dass durch die Erscheinung am 7. August
\clearpage\SourceLeaf{98a}\HistoricalHead{93}{R}{GAUSS AN ENCKE.}
\noindent \(7\frac{3}{4}\) Uhr Abends der Längenunterschied zwischen Göttingen und anderen Orten
wohl viel genauer als durch Monds oder Trabanten Finsternisse bestimmt
werden könnte. \textsc{Webers} Aufzeichnungen für \(7^{\mathrm{h}}\,40'\) gaben immer für jede
Schwingungsdauer \((20'')\) zwei Scalentheile Wachsthum.

Alle hiesigen Theilnehmer erwarten nun mit Verlangen die correspondirenden
Beobachtungen, besonders von Ihnen. Der Leipziger Apparat scheint
am 22. Junius noch nicht hinlänglich gegen den Luftzug geschützt gewesen
zu sein, hoffentlich diesmal besser. \textsc{Sartorius} und \textsc{Listing} werden in München
beobachtet haben. \mbox{. . . . . . . . . . . . . . . . . . . . .}
\clearpage\endgroup
